% ============================================================================
% SEC05.TEX - Section 3.5: Object-Oriented Programming
% ============================================================================

\section{Object-Oriented Programming dalam Unity}

\begin{learningoutcomes}
\item Memahami konsep class dan object
\item Menggunakan inheritance
\item Memahami polymorphism
\item Menerapkan encapsulation
\end{learningoutcomes}

\subsection{Class dan Object}

\begin{lstlisting}[language={[Sharp]C}]
public class Enemy : MonoBehaviour
{
    public int health = 100;
    public float speed = 2f;
    
    public void TakeDamage(int damage)
    {
        health -= damage;
        if (health <= 0)
        {
            Die();
        }
    }
    
    private void Die()
    {
        Destroy(gameObject);
    }
}
\end{lstlisting}

\subsection{Inheritance}

\begin{lstlisting}[language={[Sharp]C}]
public class BaseEnemy : MonoBehaviour
{
    protected int health;
    
    public virtual void Attack()
    {
        Debug.Log("Enemy attacks");
    }
}

public class MeleeEnemy : BaseEnemy
{
    public override void Attack()
    {
        Debug.Log("Melee attack");
    }
}
\end{lstlisting}

\begin{catatan}
Gunakan inheritance untuk membuat hierarki class yang reusable. Gunakan virtual/override untuk polymorphism.
\end{catatan}
